\cotefile{156-6}
\scanpage{2}\mspage{1}
\margin{\tikz[baseline=(g.base)]{\node[inner sep=2pt] (g) {GF \uline{VI}}; \draw (g.north west)--(g.south west)--(g.south east);}}
\begin{center}Analysis situs (deuxième mouture)\end{center}

% slanted
\margin{18 juin 86}Avant de décrire ce qu'est une \uline{contrée}, je vais décrire ce qui sera une ``contrée avec notion de multistrates'' -- la famille des multistrates \inserted{choisies} jouant un peu le rôle \unsure{de'|d'|des} une famille d'ouverts \deleted{\unsure{basiques|basique} pour} \inserted{engendrant} une topologie, ou une famille génératrice d'éléments d'un topos. Je vais donc commencer par

\textcircled{I}\quad ``Algèbre de figures'' ou ``\uline{\deleted{\unsure{a}}Atelier}''.

1.\quad Une \uline{algèbre de figures} implique avant tout trois types d'objets, les \uline{lieux}, les \uline{multistrates}, les \uline{figures}, \unsure{formant} trois ensembles
\[
\mathcal{L},\mathcal{M},\mathcal{F},\tag*{(1.1\unsure{.})}
\]
liés entre eux par diverses applications, et chacun muni de diverses structures. Ainsi, on a des applications canoniques injectives
\diagram{\begin{tikzcd}[ampersand replacement=\&,column sep=large]
{\text{(1.2)}} \& \mathcal{L} \arrow[r,hook,"\unsure{b)}"] \& \mathcal{M} \arrow[r,hook,"\unsure{a)}"] \& \mathcal{F}{,}
\end{tikzcd}}
que nous utiliserons souvent pour identifier un lieu à une multistrate particulière, et une multistrate à une figure particulière -- $\mathcal{L}$ à un \unsure{sous-ens.} de $\mathcal{M}$, $\mathcal{M}$ à un \unsure{sous-ens} de $\mathcal{F}$.

Il y a d'autre part deu\deleted{\unsure{x}}x\unsure{.} autres paires d'applications, que voici.
